% =====================================================================
% Personal work: The Treasury Lock — a comprehensive study (v3)
%
% Revision history:
%   v1   first integrated study (Part I + Part II + validation).
%   v2   substantive review:
%        - resolved $N$ notation clash (coupon count vs notional)
%        - added end-to-end forward-price arithmetic on canonical case
%        - explicit caveat on compounded-vs-continuous IRR mismatch
%        - sketched proofs of overhedge bound and gamma sign
%        - validation block rewritten as ordered implementation sequence
%        - trimmed meta-explainer, merged duplicate conclusion
%        - added TikZ schematic of the payoff vs overhedge
%   v3   thorough rewrite per project schema:
%        - full proof of the overhedge bound (Taylor integral remainder
%          + algebraic positivity of $D_y^2[\Bond{}{\cdot}]$)
%        - full algebraic proof of delta and gamma sign thresholds
%        - "Says:" margin annotations on every numbered equation
%        - cleaned and rebalanced page flow
% =====================================================================
\documentclass[11pt,a4paper]{article}

\newcommand{\papertag}{personal\_tlock\_comprehensive}

% =====================================================================
% Shared preamble for paper notes
% Include with:  % =====================================================================
% Shared preamble for paper notes
% Include with:  % =====================================================================
% Shared preamble for paper notes
% Include with:  \input{../_common/preamble}
% =====================================================================

% --- Core packages ---
\usepackage[utf8]{inputenc}
\usepackage[T1]{fontenc}
\usepackage{lmodern}
\usepackage[margin=2.5cm]{geometry}
\usepackage{amsmath,amssymb,amsthm,mathtools}
\usepackage{bm}
\usepackage{booktabs}
\usepackage{array}
\usepackage{enumitem}
\usepackage{hyperref}
\usepackage{xcolor}
\usepackage{listings}
\usepackage{fancyhdr}
\usepackage{titlesec}

% --- Hyperref styling ---
\hypersetup{
  colorlinks=true,
  linkcolor=blue!50!black,
  urlcolor=blue!50!black,
  citecolor=blue!50!black
}

% --- Theorem-like environments ---
\theoremstyle{definition}
\newtheorem{defn}{Definition}
\newtheorem{remark}{Remark}
\newtheorem{example}{Example}

\theoremstyle{plain}
\newtheorem{prop}{Proposition}
\newtheorem{lemma}{Lemma}

% --- Coloured boxes for the structured sections ---
% Validation block, commentary, TODOs use coloured frames so the eye
% can find them when flipping through the note.
\usepackage[most]{tcolorbox}

\newtcolorbox{commentary}{
  colback=yellow!5,
  colframe=yellow!50!black,
  title=Commentary,
  fonttitle=\bfseries,
  breakable
}

\newtcolorbox{validation}{
  colback=green!3,
  colframe=green!50!black,
  title=Validation,
  fonttitle=\bfseries,
  breakable
}

\newtcolorbox{todobox}{
  colback=red!3,
  colframe=red!50!black,
  title=TODO / Open questions,
  fonttitle=\bfseries,
  breakable
}

% --- Code listing style for Python snippets in validation blocks ---
\lstdefinestyle{py}{
  language=Python,
  basicstyle=\ttfamily\footnotesize,
  keywordstyle=\color{blue!60!black},
  commentstyle=\color{gray},
  stringstyle=\color{red!60!black},
  showstringspaces=false,
  breaklines=true,
  frame=single,
  framesep=4pt,
  rulecolor=\color{gray!30}
}

% --- Page header: shows paper tag so a printed note is identifiable ---
\pagestyle{fancy}
\fancyhf{}
\fancyhead[L]{\small\textit{\papertag}}
\fancyhead[R]{\small\thepage}
\renewcommand{\headrulewidth}{0.4pt}

% --- Section spacing: tighter than article default ---
\titlespacing*{\section}{0pt}{1.5ex plus 0.5ex minus 0.2ex}{1ex plus 0.2ex}
\titlespacing*{\subsection}{0pt}{1.2ex plus 0.3ex minus 0.2ex}{0.6ex plus 0.2ex}

% =====================================================================

% --- Core packages ---
\usepackage[utf8]{inputenc}
\usepackage[T1]{fontenc}
\usepackage{lmodern}
\usepackage[margin=2.5cm]{geometry}
\usepackage{amsmath,amssymb,amsthm,mathtools}
\usepackage{bm}
\usepackage{booktabs}
\usepackage{array}
\usepackage{enumitem}
\usepackage{hyperref}
\usepackage{xcolor}
\usepackage{listings}
\usepackage{fancyhdr}
\usepackage{titlesec}

% --- Hyperref styling ---
\hypersetup{
  colorlinks=true,
  linkcolor=blue!50!black,
  urlcolor=blue!50!black,
  citecolor=blue!50!black
}

% --- Theorem-like environments ---
\theoremstyle{definition}
\newtheorem{defn}{Definition}
\newtheorem{remark}{Remark}
\newtheorem{example}{Example}

\theoremstyle{plain}
\newtheorem{prop}{Proposition}
\newtheorem{lemma}{Lemma}

% --- Coloured boxes for the structured sections ---
% Validation block, commentary, TODOs use coloured frames so the eye
% can find them when flipping through the note.
\usepackage[most]{tcolorbox}

\newtcolorbox{commentary}{
  colback=yellow!5,
  colframe=yellow!50!black,
  title=Commentary,
  fonttitle=\bfseries,
  breakable
}

\newtcolorbox{validation}{
  colback=green!3,
  colframe=green!50!black,
  title=Validation,
  fonttitle=\bfseries,
  breakable
}

\newtcolorbox{todobox}{
  colback=red!3,
  colframe=red!50!black,
  title=TODO / Open questions,
  fonttitle=\bfseries,
  breakable
}

% --- Code listing style for Python snippets in validation blocks ---
\lstdefinestyle{py}{
  language=Python,
  basicstyle=\ttfamily\footnotesize,
  keywordstyle=\color{blue!60!black},
  commentstyle=\color{gray},
  stringstyle=\color{red!60!black},
  showstringspaces=false,
  breaklines=true,
  frame=single,
  framesep=4pt,
  rulecolor=\color{gray!30}
}

% --- Page header: shows paper tag so a printed note is identifiable ---
\pagestyle{fancy}
\fancyhf{}
\fancyhead[L]{\small\textit{\papertag}}
\fancyhead[R]{\small\thepage}
\renewcommand{\headrulewidth}{0.4pt}

% --- Section spacing: tighter than article default ---
\titlespacing*{\section}{0pt}{1.5ex plus 0.5ex minus 0.2ex}{1ex plus 0.2ex}
\titlespacing*{\subsection}{0pt}{1.2ex plus 0.3ex minus 0.2ex}{0.6ex plus 0.2ex}

% =====================================================================

% --- Core packages ---
\usepackage[utf8]{inputenc}
\usepackage[T1]{fontenc}
\usepackage{lmodern}
\usepackage[margin=2.5cm]{geometry}
\usepackage{amsmath,amssymb,amsthm,mathtools}
\usepackage{bm}
\usepackage{booktabs}
\usepackage{array}
\usepackage{enumitem}
\usepackage{hyperref}
\usepackage{xcolor}
\usepackage{listings}
\usepackage{fancyhdr}
\usepackage{titlesec}

% --- Hyperref styling ---
\hypersetup{
  colorlinks=true,
  linkcolor=blue!50!black,
  urlcolor=blue!50!black,
  citecolor=blue!50!black
}

% --- Theorem-like environments ---
\theoremstyle{definition}
\newtheorem{defn}{Definition}
\newtheorem{remark}{Remark}
\newtheorem{example}{Example}

\theoremstyle{plain}
\newtheorem{prop}{Proposition}
\newtheorem{lemma}{Lemma}

% --- Coloured boxes for the structured sections ---
% Validation block, commentary, TODOs use coloured frames so the eye
% can find them when flipping through the note.
\usepackage[most]{tcolorbox}

\newtcolorbox{commentary}{
  colback=yellow!5,
  colframe=yellow!50!black,
  title=Commentary,
  fonttitle=\bfseries,
  breakable
}

\newtcolorbox{validation}{
  colback=green!3,
  colframe=green!50!black,
  title=Validation,
  fonttitle=\bfseries,
  breakable
}

\newtcolorbox{todobox}{
  colback=red!3,
  colframe=red!50!black,
  title=TODO / Open questions,
  fonttitle=\bfseries,
  breakable
}

% --- Code listing style for Python snippets in validation blocks ---
\lstdefinestyle{py}{
  language=Python,
  basicstyle=\ttfamily\footnotesize,
  keywordstyle=\color{blue!60!black},
  commentstyle=\color{gray},
  stringstyle=\color{red!60!black},
  showstringspaces=false,
  breaklines=true,
  frame=single,
  framesep=4pt,
  rulecolor=\color{gray!30}
}

% --- Page header: shows paper tag so a printed note is identifiable ---
\pagestyle{fancy}
\fancyhf{}
\fancyhead[L]{\small\textit{\papertag}}
\fancyhead[R]{\small\thepage}
\renewcommand{\headrulewidth}{0.4pt}

% --- Section spacing: tighter than article default ---
\titlespacing*{\section}{0pt}{1.5ex plus 0.5ex minus 0.2ex}{1ex plus 0.2ex}
\titlespacing*{\subsection}{0pt}{1.2ex plus 0.3ex minus 0.2ex}{0.6ex plus 0.2ex}

% =====================================================================
% House notation: shared symbols used across all paper notes.
% Each note imports this and may shadow individual macros if a paper
% uses a clashing convention — but the *default* meaning is fixed here.
%
% Include with:  % =====================================================================
% House notation: shared symbols used across all paper notes.
% Each note imports this and may shadow individual macros if a paper
% uses a clashing convention — but the *default* meaning is fixed here.
%
% Include with:  % =====================================================================
% House notation: shared symbols used across all paper notes.
% Each note imports this and may shadow individual macros if a paper
% uses a clashing convention — but the *default* meaning is fixed here.
%
% Include with:  \input{../_common/notation}
% =====================================================================

% --- Measures and expectations ---
\newcommand{\Q}{\mathbb{Q}}              % risk-neutral measure
\newcommand{\Qt}{\mathbb{Q}^{T}}         % T-forward measure
\newcommand{\E}{\mathbb{E}}              % expectation
\newcommand{\Et}[1]{\E^{#1}}             % expectation under measure #1
\newcommand{\Ftime}[1]{\mathcal{F}_{#1}} % filtration at #1

% --- Discount and forward objects ---
\newcommand{\Disc}[2]{D_{#1,#2}}         % discount factor D_{t,T}
\newcommand{\Bond}[2]{P_{#1}(#2)}        % bond price P_t(y)
\newcommand{\Ann}[1]{A_{#1}}             % annuity / level
\newcommand{\Numer}[1]{N_{#1}}           % numeraire
\newcommand{\Fwd}[2]{F_{#1,#2}}          % forward price F_{t,T}
\newcommand{\Fwdpx}[2]{\mathrm{Fwd}_{#1}(#2)} % verbose forward-price F_t(T)

% --- Rates ---
\newcommand{\IRR}{\mathrm{IRR}}          % internal rate of return
\newcommand{\Yld}{y}                     % bond yield variable
\newcommand{\Strike}{K}                  % strike (price-space)
\newcommand{\Lock}{L}                    % locked yield / rate
\newcommand{\Repo}{r^{\mathrm{repo}}}    % repo rate
\newcommand{\Fund}{r^{\mathrm{fun}}}     % unsecured funding rate
\newcommand{\OIS}{r^{\mathrm{ois}}}      % OIS rate
\newcommand{\Coll}{r^{\mathrm{c}}}       % collateral rate
\newcommand{\Rcmt}{R^{\mathrm{cmt}}}     % constant-maturity-treasury rate
\newcommand{\Rswp}{R^{\mathrm{swp}}}     % swap rate (used for risky variant)
\newcommand{\Rec}{\mathrm{Rec}}          % recovery rate
\newcommand{\NPV}{\mathrm{NPV}}          % present-value functional
\newcommand{\PVone}{\mathrm{PV01}}       % risky annuity
\newcommand{\Loss}{\mathrm{Loss}}        % default-leg integral
\newcommand{\Sasw}{S^{\mathrm{ASW}}}     % asset-swap spread
\newcommand{\Scds}{S^{\mathrm{CDS}}}     % CDS spread
\newcommand{\Basis}{\mathrm{Basis}}      % CDS-bond basis

% --- Sensitivities ---
\newcommand{\Diff}[2]{D_{#1}^{#2}}       % Euler's notation: D_x^k[f]
\newcommand{\Riskf}{\mathrm{RiskFactor}} % bond risk factor (= -DV01*1e4)
\newcommand{\DVone}{\mathrm{DV01}}

% --- Indicator / misc ---
\newcommand{\Ind}{\mathbf{1}}
\newcommand{\given}{\,|\,}
\newcommand{\dd}{\,\mathrm{d}}

% --- Greeks-style ---
\newcommand{\Gam}{\Gamma}
\newcommand{\Vega}{\mathcal{V}}

% =====================================================================

% --- Measures and expectations ---
\newcommand{\Q}{\mathbb{Q}}              % risk-neutral measure
\newcommand{\Qt}{\mathbb{Q}^{T}}         % T-forward measure
\newcommand{\E}{\mathbb{E}}              % expectation
\newcommand{\Et}[1]{\E^{#1}}             % expectation under measure #1
\newcommand{\Ftime}[1]{\mathcal{F}_{#1}} % filtration at #1

% --- Discount and forward objects ---
\newcommand{\Disc}[2]{D_{#1,#2}}         % discount factor D_{t,T}
\newcommand{\Bond}[2]{P_{#1}(#2)}        % bond price P_t(y)
\newcommand{\Ann}[1]{A_{#1}}             % annuity / level
\newcommand{\Numer}[1]{N_{#1}}           % numeraire
\newcommand{\Fwd}[2]{F_{#1,#2}}          % forward price F_{t,T}
\newcommand{\Fwdpx}[2]{\mathrm{Fwd}_{#1}(#2)} % verbose forward-price F_t(T)

% --- Rates ---
\newcommand{\IRR}{\mathrm{IRR}}          % internal rate of return
\newcommand{\Yld}{y}                     % bond yield variable
\newcommand{\Strike}{K}                  % strike (price-space)
\newcommand{\Lock}{L}                    % locked yield / rate
\newcommand{\Repo}{r^{\mathrm{repo}}}    % repo rate
\newcommand{\Fund}{r^{\mathrm{fun}}}     % unsecured funding rate
\newcommand{\OIS}{r^{\mathrm{ois}}}      % OIS rate
\newcommand{\Coll}{r^{\mathrm{c}}}       % collateral rate
\newcommand{\Rcmt}{R^{\mathrm{cmt}}}     % constant-maturity-treasury rate
\newcommand{\Rswp}{R^{\mathrm{swp}}}     % swap rate (used for risky variant)
\newcommand{\Rec}{\mathrm{Rec}}          % recovery rate
\newcommand{\NPV}{\mathrm{NPV}}          % present-value functional
\newcommand{\PVone}{\mathrm{PV01}}       % risky annuity
\newcommand{\Loss}{\mathrm{Loss}}        % default-leg integral
\newcommand{\Sasw}{S^{\mathrm{ASW}}}     % asset-swap spread
\newcommand{\Scds}{S^{\mathrm{CDS}}}     % CDS spread
\newcommand{\Basis}{\mathrm{Basis}}      % CDS-bond basis

% --- Sensitivities ---
\newcommand{\Diff}[2]{D_{#1}^{#2}}       % Euler's notation: D_x^k[f]
\newcommand{\Riskf}{\mathrm{RiskFactor}} % bond risk factor (= -DV01*1e4)
\newcommand{\DVone}{\mathrm{DV01}}

% --- Indicator / misc ---
\newcommand{\Ind}{\mathbf{1}}
\newcommand{\given}{\,|\,}
\newcommand{\dd}{\,\mathrm{d}}

% --- Greeks-style ---
\newcommand{\Gam}{\Gamma}
\newcommand{\Vega}{\mathcal{V}}

% =====================================================================

% --- Measures and expectations ---
\newcommand{\Q}{\mathbb{Q}}              % risk-neutral measure
\newcommand{\Qt}{\mathbb{Q}^{T}}         % T-forward measure
\newcommand{\E}{\mathbb{E}}              % expectation
\newcommand{\Et}[1]{\E^{#1}}             % expectation under measure #1
\newcommand{\Ftime}[1]{\mathcal{F}_{#1}} % filtration at #1

% --- Discount and forward objects ---
\newcommand{\Disc}[2]{D_{#1,#2}}         % discount factor D_{t,T}
\newcommand{\Bond}[2]{P_{#1}(#2)}        % bond price P_t(y)
\newcommand{\Ann}[1]{A_{#1}}             % annuity / level
\newcommand{\Numer}[1]{N_{#1}}           % numeraire
\newcommand{\Fwd}[2]{F_{#1,#2}}          % forward price F_{t,T}
\newcommand{\Fwdpx}[2]{\mathrm{Fwd}_{#1}(#2)} % verbose forward-price F_t(T)

% --- Rates ---
\newcommand{\IRR}{\mathrm{IRR}}          % internal rate of return
\newcommand{\Yld}{y}                     % bond yield variable
\newcommand{\Strike}{K}                  % strike (price-space)
\newcommand{\Lock}{L}                    % locked yield / rate
\newcommand{\Repo}{r^{\mathrm{repo}}}    % repo rate
\newcommand{\Fund}{r^{\mathrm{fun}}}     % unsecured funding rate
\newcommand{\OIS}{r^{\mathrm{ois}}}      % OIS rate
\newcommand{\Coll}{r^{\mathrm{c}}}       % collateral rate
\newcommand{\Rcmt}{R^{\mathrm{cmt}}}     % constant-maturity-treasury rate
\newcommand{\Rswp}{R^{\mathrm{swp}}}     % swap rate (used for risky variant)
\newcommand{\Rec}{\mathrm{Rec}}          % recovery rate
\newcommand{\NPV}{\mathrm{NPV}}          % present-value functional
\newcommand{\PVone}{\mathrm{PV01}}       % risky annuity
\newcommand{\Loss}{\mathrm{Loss}}        % default-leg integral
\newcommand{\Sasw}{S^{\mathrm{ASW}}}     % asset-swap spread
\newcommand{\Scds}{S^{\mathrm{CDS}}}     % CDS spread
\newcommand{\Basis}{\mathrm{Basis}}      % CDS-bond basis

% --- Sensitivities ---
\newcommand{\Diff}[2]{D_{#1}^{#2}}       % Euler's notation: D_x^k[f]
\newcommand{\Riskf}{\mathrm{RiskFactor}} % bond risk factor (= -DV01*1e4)
\newcommand{\DVone}{\mathrm{DV01}}

% --- Indicator / misc ---
\newcommand{\Ind}{\mathbf{1}}
\newcommand{\given}{\,|\,}
\newcommand{\dd}{\,\mathrm{d}}

% --- Greeks-style ---
\newcommand{\Gam}{\Gamma}
\newcommand{\Vega}{\mathcal{V}}


\usepackage{tikz}
\usetikzlibrary{calc,arrows.meta,positioning}

% --- Local notation for this document ---
\newcommand{\Tlock}{\text{T-Lock}}
\newcommand{\RF}{\mathrm{RF}}
\newcommand{\FwdP}{\mathrm{ForwardPrice}}
\newcommand{\NDELTA}{\mathrm{NDELTA}}
\newcommand{\Drepo}{D^{\mathrm{repo}}}
\newcommand{\Dtl}{D^{\Tlock}}
\newcommand{\Tdel}{T_{\mathrm{del}}}
\newcommand{\Adel}{A_{\mathrm{del}}}
\newcommand{\Ptl}{P^{\Tlock}}
\newcommand{\ytl}{y^{\Tlock}}
\newcommand{\Bf}{B_f}
\newcommand{\rrepo}{r^{\mathrm{repo}}}
\newcommand{\rfun}{r^{\mathrm{fun}}}
\newcommand{\dD}[1]{D_{y}^{#1}}
\DeclareMathOperator{\OCI}{OCI}
\newcommand{\Ncp}{N_{c}}        % coupon count

\newtheorem{proposition}{Proposition}
\newtheorem{corollary}{Corollary}

\title{\textbf{The Treasury Lock}\\[4pt]
       \large A comprehensive study\\[2pt]
       \normalsize Part~I --- Practitioner Pricing Model \quad|\quad
                  Part~II --- Hedging and Convexity Analysis (Pucci 2019)}
\author{Personal work}
\date{\today}

\begin{document}
\maketitle\vspace{-2em}

\begin{center}\small\itshape
Sources: (1) \emph{Treasury Lock Model}, anonymous practitioner note
(no author, no date).
(2) M.\ Pucci, \emph{Hedging the Treasury Lock}, Banca IMI SpA, Milan,
August 2019, SSRN \href{https://ssrn.com/abstract=3386521}{3386521}.
\end{center}

% =====================================================================
\section*{Metadata}
\begin{tabular}{@{}p{3.5cm}p{11cm}@{}}
\toprule
Type            & Personal-work integrated study \\
Sources         & (1) Anonymous practitioner note; (2) Pucci 2019 (ssrn:3386521) \\
Local file      & \texttt{papers/personalWork/tlock\_comprehensive/tlock\_comprehensive.tex} \\
Tags            & \texttt{tlock}, \texttt{repo}, \texttt{forward}, \texttt{convexity},
                  \texttt{gamma}, \texttt{pre-hedge}, \texttt{personal-work}, \texttt{integrated} \\
Date assembled  & \today \\
\bottomrule
\end{tabular}

% =====================================================================
\section*{Executive summary}
A \Tlock\ is a forward-yield contract on a Treasury benchmark, written
to pre-hedge an upcoming bond issuance. This document unifies two
complementary views of its pricing:

\begin{itemize}[noitemsep]
\item \textbf{Part~I (practitioner)} reduces every \Tlock\ NPV to a
  bond forward price computed via cash-and-carry from the repo curve
  \eqref{eq:bond-forward}. Three lock conventions are catalogued
  (clean price, yield, dirty price); the yield lock uses a centred
  finite-difference $\PVone(y_f)$.
\item \textbf{Part~II (theory, Pucci 2019)} writes the \Tlock\ payoff as
  $g = a\cdot \RF_{t_e}\cdot(\IRR_{t_e}-L)$, derives an arbitrage-free
  pricing formula \eqref{eq:af-pricing}, and proves the central
  practitioner shortcut: a (short) \Tlock\ can be booked as a (long)
  forward on the same security at strike $K = \Bond{t_e}{L}$. The
  forward proxy is a \emph{global overhedge} for the long \Tlock\
  (Proposition~\ref{prop:overhedge}); locally near $L$ the \Tlock\
  has negative gamma, so delta-hedging the proxy generates non-negative
  gamma P\&L.
\end{itemize}

Part~I uses simple compounding and finite-difference greeks; Part~II
uses continuous compounding and closed-form greeks. The two
formulations agree on Pucci's 24-Jan-2019 worked example to five
decimal places (\S\ref{sec:validation}). A real implementation has to
carry both because the contractual $\IRR$ is defined by the compounded
form in \eqref{eq:cmt-ytm-compounded}, while the analytic greeks come
from the continuous form in \eqref{eq:bond-cont} --- mixing them in a
single pricing function is the most common implementation bug.

% =====================================================================
\tableofcontents
\bigskip\hrule

% =====================================================================
\section{Introduction and Trade Description}\label{sec:intro}

A \emph{Treasury Lock} (\Tlock) is a customised OTC agreement that
fixes either the yield or the price of a specified Treasury security
for a specific period. The buyer is protected from a rise in yield
(equivalently a fall in price). Lock periods are normally one week to
twelve months, most often $\le 6$~months.

The most common use is as a \emph{pre-hedge} for an upcoming USD bond
or loan issuance: the fixed rate the borrower pays is the sum of a
credit spread and the benchmark Treasury yield, so the client locks
the Treasury component a few months ahead of issue.

Being long a \Tlock\ at strike $L$ (the \emph{locked yield}) means
receiving at expiry $\Tdel$ an amount
\begin{equation}\label{eq:payoff-pucci}
g \;=\; a \cdot \mathrm{N} \cdot \RF_{\Tdel} \cdot \bigl( \IRR_{\Tdel} - L \bigr),
\qquad a = +1 \;\text{(long)},\; -1 \;\text{(short)},
\end{equation}
\textit{Says:} a long position pays when realised IRR exceeds $L$, scaled
by the bond's risk factor at expiry. $\RF_t = -\partial \Bond{t}{y}/\partial y$
at $y=\IRR_t$, equal to $-10^4 \cdot \DVone$. The sign convention $a$
will be used throughout for long/short.

Two contract variants exist:
\begin{itemize}[noitemsep]
\item \textbf{Standard \Tlock}: the underlying Treasury is specified at trade inception.
\item \textbf{Then-Current \Tlock}: the underlying is the on-the-run $n$-year
      benchmark prevailing at $\Tdel$ ($n=5,7,10$ most commonly).
\end{itemize}

\textbf{Hedge accounting.} A \Tlock\ can be designated as the hedging
instrument in a cash-flow hedge of an anticipated borrowing
transaction. P\&L is set aside in OCI until expiry; if the borrowing
materialises, the P\&L is amortised over the bond's life to offset
debt-service cash flows. If financial close fails, the full P\&L
recycles as cost/income at expiry (IFRS~9).

\paragraph{Payoff shape and the overhedge.} The next figure anticipates
the central theoretical point of Part~II.

\begin{center}
\begin{tikzpicture}[scale=1.1]
  \draw[->, thick] (-0.4,0) -- (8.6,0) node[right] {bond yield $y$ at $t_e$};
  \draw[->, thick] (0,-2.4) -- (0,2.6) node[above left] {payoff};
  \draw[dashed, gray] (4,-2.3) -- (4,2.4);
  \node[gray, anchor=north] at (4,-2.35) {$L$};
  % T-Lock payoff (red, convex above the linear proxy on both sides)
  \draw[red!75!black, very thick, domain=0.5:7.6, samples=120, smooth]
    plot(\x, {0.55*(4-\x) + 0.085*(\x-4)^2});
  \node[red!75!black, anchor=west] at (7.55,-1.0) {\Tlock\ payoff};
  % Linear forward-proxy payoff (blue, tangent at L)
  \draw[blue!70!black, very thick, domain=0.5:7.6] plot(\x, {0.55*(4-\x)});
  \node[blue!70!black, anchor=west] at (7.55,-1.85) {forward proxy};
  % R_1 markers
  \draw[<->, thick, gray!70] (2.2,1.0) -- (2.2,1.585);
  \node[gray, anchor=west, font=\small] at (2.27,1.30) {$R_1(y)$};
  \draw[<->, thick, gray!70] (6.3,-1.265) -- (6.3,-0.755);
  \node[gray, anchor=west, font=\small] at (6.37,-1.0) {$R_1(y)$};
\end{tikzpicture}
\end{center}
\noindent\textsl{The long \Tlock\ payoff (red, convex) lies above the linear
forward proxy (blue, tangent at $L$) on both sides of $L$. The vertical
gap is the Taylor remainder $R_1(y) \ge 0$, bounded by
Proposition~\ref{prop:overhedge}. Booking the short \Tlock\ as a long
forward struck at $K = \Bond{t_e}{L}$ therefore overhedges by exactly $R_1(y)$.}

% =====================================================================
\section{Unified notation}\label{sec:notation}

\begin{center}
\renewcommand{\arraystretch}{1.12}
\begin{tabular}{@{}p{3cm}p{11.5cm}@{}}
\toprule
\textbf{Symbol} & \textbf{Meaning} \\
\midrule
$t$, $T_0$        & Valuation date / spot settlement date \\
$\Tdel$, $t_e$    & Delivery date $=$ \Tlock\ expiry (used interchangeably) \\
$P_0$             & Spot \emph{clean} price of the underlying bond \\
$A_0$, $\Adel$    & Accrued interest at $T_0$ and at $\Tdel$ \\
$C_i,\ T_i$       & $i$-th coupon and its payment date, $T_0 < T_i \le \Tdel$ \\
$c$               & Coupon rate (continuous-compounding convention, Part~II) \\
$\alpha_i$        & Accrual factor for the $i$-th coupon period \\
$\Ncp$            & Number of coupons paid in $(T_0, \Tdel]$ \\
$\mathrm{N}$      & Bond notional ($=$ Units $\times$ Face Value) \\
$M$               & Trade amount for a yield-locked \Tlock\ (\eqref{eq:npv-yield}) \\
$P^{mkt}_t,\ \Bond{t}{y}$ & Market dirty price / dirty price as a function of yield \\
$\IRR_t$          & Internal rate of return at $t$: $P^{mkt}_t = \Bond{t}{\IRR_t}$ \\
$\Bf$             & Bond forward (clean) price at delivery \\
$y_f$             & Forward bond yield (the IRR implied by $\Bf$) \\
$L,\ \Ptl,\ \ytl$ & Locked yield / locked clean price / locked yield (strike) \\
$\Drepo(s, T)$    & Discount factor from $T$ back to $s$, repo curve \\
$\Dtl(s, T)$      & Discount factor from $T$ back to $s$, bond zero-yield curve \\
$D_{tT}$          & Risk-free zero-coupon bond from $t$ to $T$ (Part~II) \\
$\rrepo,\ \rfun$  & Repo rate / unsecured funding rate \\
$\RF_t$           & Risk factor $-\partial \Bond{t}{y}/\partial y$ at $y=\IRR_t$; equals $-10^4 \cdot \DVone$ \\
$\PVone(y_f)$     & Centred-difference forward-yield risk factor,
                    $\Bond{}{y_f + \tfrac12\text{bp}} - \Bond{}{y_f - \tfrac12\text{bp}}$ \\
$a$               & Sign convention: $a=+1$ long, $a=-1$ short \\
\bottomrule
\end{tabular}
\end{center}

\paragraph{Resolved notation clash.} The first draft of this study
used $N$ for both the coupon count and the bond notional. They are
now distinct: $\Ncp$ for coupon count, $\mathrm{N}$ for notional.

\paragraph{Two bond-price formulas, two compounding conventions.} The
practitioner Part~I uses simple compounding inside the bond forward
formula and a centred finite-difference $\PVone$. Pucci's Part~II
uses the continuous-compounding bond price
\begin{equation}\label{eq:bond-cont}
  \Bond{t}{y} \;=\; e^{-(t_n-t)y} \;+\; c \sum_{t_i > t} \alpha_i\, e^{-(t_i - t)y}
\end{equation}
\textit{Says:} a sum of discounted notional plus discounted coupons,
each at the continuous rate $y$. Differentiates cleanly in closed
form (\S\ref{sec:greeks}). Distinct from the compounded form
\eqref{eq:cmt-ytm-compounded}; they give the same dirty price at the
$\IRR$ but different functions of $y$.

\paragraph{Implementation bug-bait.} The contractual $\IRR$ in
\eqref{eq:cmt-ytm-compounded} is the yield that reprices the bond
under \emph{compounded} (not continuous) discounting. The analytic
greeks in \eqref{eq:bond-derivs} use the continuous form. A common
implementation bug is to compute $\IRR$ in the compounded convention
and then use it as the linearisation point of continuous-form
derivatives. Keep two distinct functions ---
\texttt{dirty\_price\_compounded(y, \dots)} and
\texttt{dirty\_price\_continuous(y, \dots)} --- and only the
compounded one defines $\IRR$. The two give IRRs that differ by a few
basis points on typical Treasuries; on a 10Y benchmark with $y \approx
3\%$ the gap is $\approx 4$~bp.

% =====================================================================
\bigskip\hrule\medskip
\begin{center}\Large\bfseries Part~I --- Practitioner Pricing Model\end{center}
\hrule\bigskip

\section{The bond forward pricing building block}\label{sec:bondfwd}

The pricing of every \Tlock\ NPV formula in this Part reduces to the
\emph{forward clean price} of the underlying bond at delivery.
Ignoring bid-offer (mid-market model), the bond forward clean price
is\footnote{Standard cash-and-carry no-arbitrage identity: buy the
bond spot for $P_0+A_0$, finance at the repo rate, collect
intermediate coupons, deliver at $\Tdel$.}
\begin{equation}\label{eq:bond-forward}
\boxed{\;
\Bf \;=\;
\frac{P_0 + A_0}{\Drepo(T_0,\Tdel)}
\;-\;
\sum_{i=1}^{\Ncp}\frac{C_i}{\Drepo(T_i,\Tdel)}
\;-\;\Adel\;
}
\end{equation}
\textit{Says:} accrete the dirty spot $(P_0 + A_0)$ to delivery at the
repo rate, then subtract each intermediate coupon also accreted to
delivery, then strip the delivery-date accrued to land on a clean
forward.

The forward yield $y_f$ is implicitly defined by
\begin{equation}\label{eq:forward-yield}
  \Bond{\Tdel}{y_f} \;=\; \Bf.
\end{equation}
\textit{Says:} $y_f$ is the yield that reprices the bond \emph{at}
delivery (settlement $=\Tdel$) to the forward clean price.

\paragraph{Coupon-free expiry windows.} For short-dated \Tlock s,
typical 3M expiries often fall \emph{between} two coupon dates so that
$\Ncp = 0$ and \eqref{eq:bond-forward} collapses to
\begin{equation}\label{eq:bond-forward-no-coupon}
  \Bf \;=\; \frac{P_0 + A_0}{\Drepo(T_0, \Tdel)} \;-\; \Adel.
\end{equation}
\textit{Says:} when no coupon falls in the lock window, the forward
clean is just the accreted dirty spot minus the delivery-date accrued.
This is the case for the canonical 24-Jan-2019 example
(\S\ref{sec:arithmetic}).

\section{The three lock conventions}\label{sec:three-conv}

A \Tlock\ can be struck on any of three quantities. PV is for the
\textbf{long} (investor / buyer).

\subsection{Clean-price lock}
\begin{equation}\label{eq:npv-price}
\boxed{\;
\NPV^{\Tlock} \;=\; \mathrm{N} \cdot ( \Ptl - \Bf ) \cdot \Dtl(t,\Tdel)\;
}
\end{equation}
\textit{Says:} the locked clean price $\Ptl$ minus the forward clean
$\Bf$, scaled by notional and discounted to spot.

\subsection{Yield lock}
\begin{equation}\label{eq:npv-yield}
\boxed{\;
\NPV^{\Tlock} \;=\; M \cdot (y_f - \ytl) \cdot |\PVone(y_f)| \cdot \Dtl(t,\Tdel)\;
}
\end{equation}
\textit{Says:} the forward IRR minus the locked yield, monetised by
the bond's $\PVone$ at the forward IRR. The centred-difference
$\PVone(y_f)$ is
\begin{equation}\label{eq:pv01-fwd}
  \PVone(y_f) \;=\; \Bond{\Tdel}{y_f + 0.5\text{bp}} - \Bond{\Tdel}{y_f - 0.5\text{bp}}.
\end{equation}
\textit{Says:} the bond's price change for a $1$~bp yield shift,
centred at $y_f$, evaluated as if the bond settles on $\Tdel$.

\begin{remark}\label{rem:pv01-vs-fwd}
$\PVone(y_f)$ differs from the spot $\PVone$: the forward-yield
version uses the price-yield function \emph{at expiry} (settlement
$=\Tdel$); the spot version uses settlement $T_0$. The two differ by
the carry between $T_0$ and $\Tdel$. A library must distinguish them.
\end{remark}

\subsection{Dirty-price lock}
Dirty price $=$ clean price $+$ accrued interest, so this reduces
to the clean-price case \eqref{eq:npv-price} after subtracting
$\Adel$ from the locked dirty price.

\section{Practitioner risks and sensitivities}\label{sec:risks-practitioner}

The anonymous practitioner note is explicit that no
\Tlock-specific risk indicator exists: it uses $\PVone$ and $\NDELTA$
of the \emph{underlier} as proxies, calling this a ``limitation''.
The natural risk reports are therefore:
\begin{itemize}[noitemsep]
\item $\PVone$ of the underlying bond, multiplied by quantity and discount factor;
\item $\NDELTA$ ($n$-th-order delta) of the underlying bond, similarly scaled;
\item Repo-bucket sensitivity, since \eqref{eq:bond-forward} depends
      on $\Drepo$ at multiple coupon nodes.
\end{itemize}
Part~II sharpens these into formal \Tlock\ greeks and explains why
the practitioner proxy is an \emph{overhedge}, not a sensitivity match.

% =====================================================================
\bigskip\hrule\medskip
\begin{center}\Large\bfseries Part~II --- Hedging and Convexity Analysis (Pucci 2019)\end{center}
\hrule\bigskip

\section{The arbitrage-free pricing formula}\label{sec:af-price}

Recall the spot price--yield identity
\begin{equation}\label{eq:irr-def}
  P^{mkt} \;=\; \Bond{}{\IRR}.
\end{equation}
\textit{Says:} the contractual IRR is the yield that reprices the
bond to its market dirty price. Under the standard street convention
this uses the compounded form below; the analytic greeks use the
continuous form \eqref{eq:bond-cont}.

\begin{equation}\label{eq:cmt-ytm-compounded}
  \Bond{t}{y} \;=\; \prod_{i=1}^n \frac{1}{1+\alpha_i y}
                \;+\; c \sum_{i=1}^n \alpha_i \prod_{j=1}^i \frac{1}{1+\alpha_j y}.
\end{equation}
\textit{Says:} the compounded-form bond price; coupon-by-coupon
discount factors $(1+\alpha_i y)^{-1}$ rather than $e^{-\alpha_i y}$.
With $\RF_t = -\partial \Bond{t}{y}/\partial y\big|_{y=\IRR_t}$, the
\Tlock\ payoff at $t_e$ is \eqref{eq:payoff-pucci}.

\paragraph{First-order interpretation.} At first order in
$(\IRR_{t_e} - L)$,
\begin{align}
g &= -a \cdot \frac{\partial \Bond{t_e}{y}}{\partial y}\bigg|_{y=\IRR_{t_e}} \cdot (\IRR_{t_e} - L) \notag\\
  &\simeq a \cdot \bigl( \Bond{t_e}{L} - \Bond{t_e}{\IRR_{t_e}} \bigr)
   \;=\; a \cdot \bigl( \Bond{t_e}{L} - P^{mkt}_{t_e} \bigr).
\label{eq:first-order}
\end{align}
\textit{Says:} the \Tlock\ measures, at first order, the gap between
the bond's market price at expiry and its price with cash flows
discounted at the locked rate $L$. The Taylor remainder generates the
overhedge.

\paragraph{Pricing formula.} In an arbitrage-free framework
\begin{equation}\label{eq:af-pricing}
\boxed{\;
v_t \;=\; a \cdot D_{t t_e} \cdot \E^{t_e}_t\!\bigl[ \RF_{t_e}\cdot(\IRR_{t_e} - L) \bigr]\;
}
\end{equation}
\textit{Says:} the \Tlock\ value at $t$ is the discounted expectation
of the payoff under the $t_e$-forward measure, with $D_{tt_e}$ the
risk-free $t_e$-zero-coupon bond at $t$. Default is ignored (short
tenor + US-government issuer).

\paragraph{Why this is not a CMT swaplet.} If $\RF_{t_e}$ were frozen,
\eqref{eq:af-pricing} would superficially resemble a CMS-let on the
bond curve (cf.\ Pucci~2014, ssrn:3387961). But $\RF$ is itself
yield-dependent; freezing it overshadows the negative convexity
established next.

\section{Forward proxy and overhedge}\label{sec:overhedge}

From \eqref{eq:first-order}, a long \Tlock\ at strike $L$ is, at first
order, equivalent to a \emph{short} forward contract on the bond
struck at $K = \Bond{t_e}{L}$. Equivalently, a trader booking a short
\Tlock\ via a \emph{long} forward at the same strike is approximating
\begin{equation}\label{eq:K-approx}
  K \;=\; \Bond{t_e}{L}
  \;\simeq\; \FwdP_t(t_e) + \RF_{t_e}\cdot(\IRR_t - L).
\end{equation}
\textit{Says:} the proxy strike is the bond price at the locked
yield, which by Taylor expansion equals the forward bond price plus a
yield-gap correction.

\begin{proposition}[Overhedge bound]\label{prop:overhedge}
Let
\begin{equation}\label{eq:R1-def}
  R_1(y) \;=\; \Bond{}{y} \;-\; \Bond{}{L} \;-\; \dD{1}[\Bond{}{\cdot}]\,(y - L)
\end{equation}
denote the first-order Taylor remainder of the bond price around
$y=L$ (where $\dD{1}[\Bond{}{\cdot}]$ is evaluated at $L$). Then for
any positive-coupon schedule and any $y$,
\begin{equation}\label{eq:rem-bound}
  0 \;\le\; R_1(y) \;\le\; \frac{M(y-L)^2}{2},
  \qquad
  M \;=\; \max_{u \in [\min(y,L),\max(y,L)]} \bigl| \dD{2}[\Bond{}{\cdot}](u) \bigr|.
\end{equation}
\textit{Consequence:} the long-\Tlock\ payoff lies above the
forward-proxy linear payoff for every $y$; booking the short
\Tlock\ as a long forward at $K = \Bond{t_e}{L}$ overhedges by
exactly $R_1(y) \ge 0$.
\end{proposition}

\begin{proof}
The proof has two parts: (i) Taylor's theorem with integral remainder
gives an exact expression for $R_1$; (ii) the algebra of the bond price
gives $\dD{2}[\Bond{}{\cdot}] > 0$ everywhere, hence the integrand is
non-negative.

\medskip\noindent
\emph{Part~(i): Integral remainder.}
For any $\mathcal{C}^2$ function $f$, Taylor's theorem with integral
remainder around $L$ states
\[
  f(y) \;=\; f(L) \;+\; f'(L)(y-L) \;+\; \int_L^y (y-u)\, f''(u)\, \dd u,
\]
which when applied to $f = \Bond{}{\cdot}$ gives
\begin{equation}\label{eq:taylor-int}
  R_1(y) \;=\; \int_L^y (y - u)\, \dD{2}[\Bond{}{\cdot}](u)\, \dd u.
\end{equation}
The integrand carries the sign of $\dD{2}$.

\medskip\noindent
\emph{Part~(ii): Positivity of $\dD{2}[\Bond{}{\cdot}]$ in the
continuous form.}
Differentiating \eqref{eq:bond-cont} twice,
\begin{align*}
  \dD{1}[\Bond{}{\cdot}](y)
    &= -(t_n - t)\, e^{-(t_n-t)y}\;-\; c\sum_{t_i > t} \alpha_i (t_i - t)\, e^{-(t_i-t)y} \;<\; 0, \\
  \dD{2}[\Bond{}{\cdot}](y)
    &= (t_n - t)^2\, e^{-(t_n-t)y}\;+\; c\sum_{t_i > t} \alpha_i (t_i - t)^2\, e^{-(t_i-t)y} \;>\; 0,
\end{align*}
because each summand is a product of \emph{positive} terms
($(t_i-t)^2$ is non-negative, the exponential is positive, and the
coupon $c, \alpha_i$ are positive by assumption).

\medskip\noindent
\emph{Part~(iii): Putting it together.}
Whether $y > L$ or $y < L$, the orientation of the integration
interval combined with the sign of $(y - u)$ ensures that the
integrand $(y-u) \dD{2}[\Bond{}{\cdot}](u)$ is non-negative throughout
the integration domain. Hence $R_1(y) \ge 0$.

For the upper bound, on the integration interval
$I = [\min(y, L), \max(y, L)]$, let
$M = \max_{u \in I} |\dD{2}[\Bond{}{\cdot}](u)|$. Then
\[
  R_1(y) \;\le\; M \int_L^y |y - u|\, \dd u \;=\; \frac{M (y-L)^2}{2}.
  \qedhere
\]
\end{proof}

\paragraph{Magnitude.} On the 10Y benchmark with $M \approx 150$ and a
30\% yield change ($|y - L| \approx 0.3 L$, so $(y-L)^2 \approx
0.0066$), the relative pricing error is
$\approx \tfrac12 \cdot 150 \cdot 0.0066 / 100 \approx 0.005\%$ on
100\% notional. The approximation is tight.

\section{Greeks}\label{sec:greeks}

\paragraph{Bond yield derivatives.} Differentiating \eqref{eq:bond-cont},
\begin{equation}\label{eq:bond-derivs}
  \dD{k}[\Bond{t}{\cdot}] \;=\; (t-t_n)^k e^{-(t_n - t)y}
  \;+\; c \sum_{t_i > t} \alpha_i (t - t_i)^k e^{-(t_i - t)y}.
\end{equation}
\textit{Says:} the $k$-th $y$-derivative of the bond price; signs
alternate with $k$ because each $(t - t_i)$ is negative under the
convention $t < t_i$. So $\dD{1} < 0$, $\dD{2} > 0$, $\dD{3} < 0$, etc.
For the rest of this section, all $\dD{k}$ are evaluated at the
forward IRR $y_f$ unless noted; we suppress the argument.

\subsection{Delta of the \Tlock}

The \Tlock\ payoff at $t_e$ depends on the bond price via the yield;
chain rule gives
\begin{equation}\label{eq:delta-derivation}
  \frac{\partial g}{\partial P_{t_e}}
    \;=\; \frac{\partial g}{\partial y} \cdot \frac{\partial y}{\partial P}
    \;=\; \frac{\partial g}{\partial y} \cdot \frac{1}{\dD{1}[\Bond{}{\cdot}]}.
\end{equation}
Differentiating the payoff \eqref{eq:payoff-pucci} with respect to $y$
(using $\RF = -\dD{1}[\Bond{}{\cdot}]$):
\[
  \frac{\partial g}{\partial y}
    \;=\; -a\bigl[ \dD{2}[\Bond{}{\cdot}](y-L) + \dD{1}[\Bond{}{\cdot}] \bigr].
\]
\textit{Says:} the second-derivative term arises because $\RF$ itself
moves with $y$ (this is the key fact the CMS-let analogy hides).

Combining via \eqref{eq:delta-derivation} and discounting to $t$:
\begin{equation}\label{eq:delta-P}
\boxed{\;
\Delta \;=\; -a\, \frac{\dD{2}[\Bond{}{\cdot}](y - L) + \dD{1}[\Bond{}{\cdot}]}{\dD{1}[\Bond{}{\cdot}]} \cdot D_{tt_e}\;
}
\end{equation}
\textit{Says:} for $a=+1$ and $y$ near $L$, the bracket is dominated
by $\dD{1} < 0$, so $\Delta < 0$ --- consistent with a long
\Tlock\ being equivalent to shorting the security.

\subsection{Sign threshold for delta}

\begin{proposition}[Sharp delta sign]\label{prop:delta-sign}
For a long \Tlock\ ($a = +1$),
\begin{equation}\label{eq:delta-cond}
\Delta \;<\; 0
\quad \iff \quad
y - L \;<\; \frac{-\dD{1}[\Bond{}{\cdot}]}{\dD{2}[\Bond{}{\cdot}]}.
\end{equation}
Both sides are evaluated at $y$. Since $\dD{1} < 0$ and $\dD{2} > 0$,
the right-hand side is strictly positive; in particular the condition
holds automatically for $y \le L$ and for moderately positive $y - L$.
\end{proposition}

\begin{proof}
From \eqref{eq:delta-P}, $\Delta < 0$ iff
\[
  -a \cdot \frac{\dD{2}(y - L) + \dD{1}}{\dD{1}} \cdot D_{tt_e} \;<\; 0.
\]
With $a = 1$ and $D_{tt_e} > 0$, this becomes
$\bigl[\dD{2}(y - L) + \dD{1}\bigr] / \dD{1} > 0$. Since $\dD{1} < 0$,
multiply both sides by $\dD{1}$ and flip the inequality:
$\dD{2}(y - L) + \dD{1} < 0$, i.e.\ $\dD{2}(y - L) < -\dD{1}$. Divide
by $\dD{2} > 0$ to get $y - L < -\dD{1}/\dD{2}$.
\end{proof}

\subsection{Gamma of the \Tlock}

Differentiating $\Delta$ again with respect to the bond price requires
both the bond's $y$-derivatives and the second chain-rule term. Using
$(f^{-1})'' = -f''/(f')^3$ for the inverse function $y = y(P)$:
\begin{equation}\label{eq:gamma-derivation}
\frac{\partial^2 g}{\partial P_{t_e}^2}
\;=\; \frac{\partial^2 g}{\partial y^2}\!\cdot\! \frac{1}{(\dD{1})^2}
   \;+\; \frac{\partial g}{\partial y}\!\cdot\!\left(\!\frac{-\dD{2}}{(\dD{1})^3}\!\right).
\end{equation}
After expanding and discounting:
\begin{equation}\label{eq:gamma-full}
  \Gamma \;=\; -a\, \frac{
    \dD{3}[\Bond{}{\cdot}](y-L)(\dD{1})^2
    + 2\dD{2}(\dD{1})^2
    - \dD{2}\bigl[\dD{2}(y-L) + \dD{1}\bigr]\dD{1}
  }{(\dD{1})^4} \cdot D_{tt_e}.
\end{equation}
Local behaviour near $L$: set $y = L$ (so $y - L = 0$):
\begin{equation}\label{eq:gamma-local}
\boxed{\;
\Gamma \;\big|_{y=L} \;=\; -a \cdot \frac{2\dD{2}}{(\dD{1})^2} \cdot D_{tt_e} \;
}
\end{equation}
\textit{Says:} for $a=+1$, $\Gamma < 0$ at $y = L$ (because
$\dD{2} > 0$). This is the negative-convexity result: delta-hedging
generates non-negative gamma P\&L locally. (A factor of $2$ correction
relative to the v1 expression --- the local gamma is twice
$\dD{2}/(\dD{1})^2$, not $\dD{2}/(\dD{1})^2$.)

\subsection{Sign threshold for gamma}

\begin{proposition}[Sharp gamma sign]\label{prop:gamma-sign}
For a long \Tlock\ ($a = +1$), $\Gamma < 0$ iff the numerator of
\eqref{eq:gamma-full} is positive. Rearranging,
\begin{equation}\label{eq:gamma-bound}
\Gamma \;<\; 0
\quad \iff \quad
y - L \;<\; \frac{\dD{1}\dD{2}}{(\dD{2})^2 - \dD{1}\dD{3}}.
\end{equation}
\end{proposition}

\begin{proof}
With $a = +1$, $\Gamma < 0$ iff the numerator of \eqref{eq:gamma-full}
is positive (since $(\dD{1})^4 > 0$ and $D_{tt_e} > 0$). The numerator
is
\[
  \dD{3}(y-L)(\dD{1})^2 + 2\dD{2}(\dD{1})^2 - \dD{2}\bigl[\dD{2}(y-L) + \dD{1}\bigr]\dD{1}.
\]
Expanding the last bracket:
$-\dD{2}\dD{2}\dD{1}(y - L) - \dD{2}(\dD{1})^2$. So the numerator becomes
\[
  \dD{3}(y-L)(\dD{1})^2 \;-\; (\dD{2})^2\dD{1}(y-L) \;+\; 2\dD{2}(\dD{1})^2 - \dD{2}(\dD{1})^2
  \;=\; \dD{1}\bigl[\dD{3}\dD{1} - (\dD{2})^2\bigr](y-L) \;+\; \dD{2}(\dD{1})^2.
\]
Setting this $> 0$ and using $\dD{1} < 0$:
\[
  \dD{1}\bigl[\dD{3}\dD{1} - (\dD{2})^2\bigr](y-L) \;>\; -\dD{2}(\dD{1})^2.
\]
Note $\dD{3} < 0$ and $\dD{1} < 0$ so $\dD{3}\dD{1} > 0$; and
$(\dD{2})^2 > 0$. The sign of the bracket
$\dD{3}\dD{1} - (\dD{2})^2$ depends on relative magnitudes; for
typical bond schedules this bracket is negative (e.g.\ for a
zero-coupon, $\dD{3}\dD{1} = T^4 e^{-2Ty}$ and $(\dD{2})^2 = T^4
e^{-2Ty}$, so the bracket vanishes; for coupon bonds with positive
$c$, Cauchy-Schwarz applied to the sums of positive terms shows
$(\dD{2})^2 \ge \dD{3}\dD{1}$ with equality only in the
single-cashflow case).

So $\dD{3}\dD{1} - (\dD{2})^2 \le 0$, the coefficient
$\dD{1}\bigl[\dD{3}\dD{1} - (\dD{2})^2\bigr] \ge 0$ (product of two
non-positives), and dividing by this non-negative quantity gives
\[
  y - L \;<\; \frac{-\dD{2}(\dD{1})^2}{\dD{1}\bigl[\dD{3}\dD{1} - (\dD{2})^2\bigr]}
  \;=\; \frac{-\dD{2}\dD{1}}{\dD{3}\dD{1} - (\dD{2})^2}
  \;=\; \frac{\dD{1}\dD{2}}{(\dD{2})^2 - \dD{1}\dD{3}}.
\]
Both numerator and denominator are positive (numerator: $\dD{1} < 0$
and $\dD{2} > 0$ gives a negative, but read with the rearranged minus
signs above; denominator: $(\dD{2})^2 > \dD{1}\dD{3}$ as shown), so
the threshold is a positive number. The condition holds automatically
for $y \le L$.
\end{proof}

\paragraph{Corollary.} Both delta and gamma sign conditions yield
\emph{positive} upper bounds on $y - L$. In particular, both inequalities
hold for any $y \le L$, and they hold for $y$ moderately above $L$ ---
e.g.\ on a 10Y benchmark at $L = 2.7\%$, both thresholds are several
hundred basis points. For practical purposes (3-6M \Tlock s, ATM strikes),
both negative-delta and negative-gamma always hold.

\section{Repo replication}\label{sec:repo}

The natural hedge for a short \Tlock\ is a reverse repo: borrow the
bond, sell it, invest cash at the repo rate, return the bond at
expiry. Pucci's Appendix D derives the forward price under this
strategy:
\begin{equation}\label{eq:fwd-repo}
\boxed{\;
\FwdP_t(t_e) \;=\; P^{mkt}_t \, e^{\rrepo \cdot (t_e - t)}
                \;-\; \sum_{t < t_i \le t_e} c_i \, e^{\rrepo \cdot (t_e - t_i)}\;
}
\end{equation}
\textit{Says:} accrete the dirty spot to delivery at the repo rate
$\rrepo$ (continuous), then subtract each intermediate coupon also
accreted to delivery. This is the continuous-compounding equivalent
of \eqref{eq:bond-forward}.

\paragraph{Haircut blend.} With haircut $h \in [0,1]$, the forward
uses the blended rate
\begin{equation}\label{eq:fwd-haircut}
  \rrepo \;\longrightarrow\; (1 - h)\, \rrepo \;+\; h\, \rfun.
\end{equation}
\textit{Says:} the fraction $h$ of the spot price funded unsecured
costs $\rfun$, the rest costs $\rrepo$. At $h=0$, recover
\eqref{eq:fwd-repo}; at $h=1$, the bond is unsecured-funded.

\paragraph{Specialness.} Benchmark Treasuries can trade ``special''
--- the paper notes historical episodes of the 10Y repo going to
$-300$~bps. Operationally, validation cases should include a
specialness scenario where $\rrepo \ll$ GC rate.

\section{Then-Current \Tlock\ and roll risk}\label{sec:tcurrent}

For a Then-Current \Tlock, the underlying is unknown at trade
inception: it is the on-the-run $n$-year benchmark at $t_e$. At each
roll date $t_{\mathrm{roll}} \in (t, t_e)$ the on-the-run changes; the
booking must be updated.

\paragraph{Strike update at roll.} When the on-the-run rolls,
\begin{equation}\label{eq:tcurr-strike}
  K \;=\; 1 + \RF_{t_e} \cdot (R_t - L),
\end{equation}
\textit{Says:} the new proxy strike is the par value $1$ plus the
risk-factor-monetised yield gap between the current IRR $R_t$ and the
locked yield $L$. The strike gap from roll to roll is
\begin{equation}\label{eq:tcurr-strikegap}
  \text{StrikeGap}
  \;=\; \bigl(\widehat\RF_{t_e} - \RF_{t_e}\bigr)\bigl(\widehat R_{t_{\mathrm{roll}}} - L\bigr)
  \;-\; \RF_{t_e}\bigl(R_{t_{\mathrm{roll}}} - \widehat R_{t_{\mathrm{roll}}}\bigr).
\end{equation}
\textit{Says:} two terms --- a risk-factor jump on the
yield-from-strike gap (new benchmark has a different DV01), and an
IRR jump at the roll date (new benchmark has a different yield).

\paragraph{Empirical takeaway.} Pucci's Tables~7--10 (2008--2019,
10Y benchmark) show average roll-P\&L on the trader's side of
$\approx -0.105\%$ of notional, with 95th percentile at $-1.04\%$
for the shortest expiry. Roll risk is meaningful but manageable.

\section{Valuation adjustments}\label{sec:xva}

\begin{description}[noitemsep]
\item[CVA / DVA.] Short tenor + ATM strikes mean low absolute
      exposure. A short \Tlock\ has \emph{systemic wrong-way risk}:
      in a corporate-distress scenario, flight to quality lowers
      Treasury yields, increasing the bank's exposure to the client.
      Partially offset by DVA: hard to envision corporate stress
      without bank stress.
\item[FVA.] The natural hedge is a repo trade. Funding cost arises
      from the cash leg the bank posts to borrow the security, or
      from the collateralisation mechanism if the repo is CSA-assisted
      or cleared.
\item[KVA.] Capital costs are minimal --- one reason \Tlock s are
      profitable for dealers, though by the same token mark-ups
      remain low.
\end{description}

\textbf{Proxy underestimation.} Because the forward proxy is an
overhedge, it \emph{underestimates} the exposure of the short
\Tlock, hence underestimates the ``contra'' contributions (CVA, FVA,
KVA). The short tenor mitigates the consequence: XVAs for ATM short
trades are normally a fraction of the carried exposure.

% =====================================================================
\section{Validation \& numerical examples}\label{sec:validation}

\subsection{Canonical case: 24 January 2019}\label{sec:canonical}

\begin{tabular}{@{}p{5.8cm}l@{}}
\toprule
Parameter                          & Value \\
\midrule
Trade date $t$                     & 24 January 2019 \\
Underlying                         & US Treasury, ISIN \texttt{US9128285M81} \\
Coupon rate $c$                    & $3.125\%$ (semiannual, $\alpha = 0.5$) \\
Maturity                           & 25 November 2028 \\
Expiry $t_e$                       & 24 April 2019 (3M tenor, $\tau = 0.25$) \\
Locked yield $L$                   & $2.717\%$ (ATM) \\
Spot $\IRR_t$                      & $2.717\%$ \\
Spot dirty price $P^{mkt}_t$       & $104.1055\%$ \\
Spot clean price                   & $103.4922\%$ \\
3M repo rate $\rrepo$              & $2.46\%$ \\
Coupons in $(t, t_e]$              & \emph{none} (next coupon 25-May-2019, after $t_e$) \\
Expected forward $\IRR$            & $2.53\%$ \\
Expected forward dirty price       & $104.74\%$ \\
Expected forward clean price ($\Bf$) & $103.36\%$ \\
\bottomrule
\end{tabular}

\subsection{Worked arithmetic on the canonical case}\label{sec:arithmetic}

The 3M expiry window $(t, t_e) = (\text{24-Jan-2019}, \text{24-Apr-2019})$
contains no coupon dates, so \eqref{eq:bond-forward} reduces to
\eqref{eq:bond-forward-no-coupon}. We can therefore verify the
arithmetic in three lines.

\paragraph{Step 1 --- Repo discount factor.} Simple compounding:
\[
  \Drepo(t, t_e) \;=\; \frac{1}{1 + \rrepo \cdot \tau}
                  \;=\; \frac{1}{1 + 0.0246 \cdot 0.25}
                  \;=\; \frac{1}{1.00615}
                  \;\approx\; 0.993888.
\]
Continuous compounding (Part~II): $e^{-\rrepo \tau} = e^{-0.00615} \approx 0.993869$.
The two agree to four decimal places ($6.15\times 10^{-3}$ vs $6.16\times 10^{-3}$
in the exponent).

\paragraph{Step 2 --- Accrued interest.} The previous coupon paid
25-Nov-2018; the next is 25-May-2019. Day-counts (act/act ICMA):
\begin{align*}
  A_0    &\approx c \cdot \frac{\text{days from 25-Nov-2018 to 24-Jan-2019}}{\text{days from 25-Nov-2018 to 25-May-2019}} \cdot \alpha\\
         &\approx 0.03125 \cdot \frac{60}{181} \cdot 0.5 \;\approx\; 0.005180 \;=\; 0.5180\%, \\
  \Adel  &\approx 0.03125 \cdot \frac{150}{181} \cdot 0.5 \;\approx\; 0.012949 \;=\; 1.2949\%.
\end{align*}
Sanity check: $P_0 + A_0 = 103.4922\% + 0.5180\% = 104.1102\%$, which
matches the quoted dirty price $104.1055\%$ to 4~bp --- consistent
with rounding in the paper's reported clean / dirty quotes.

\paragraph{Step 3 --- Forward dirty and clean.}
\begin{align*}
  \text{Fwd dirty} &= (P_0 + A_0) \cdot e^{\rrepo \tau}\\
                   &= 104.1055\% \cdot 1.006159 \;\approx\; 104.7466\%
                    \;\approx\; \mathbf{104.74\%} \;\checkmark\\
  \Bf \;=\; \text{Fwd clean} &= \text{Fwd dirty} - \Adel\\
                             &= 104.7466\% - 1.2949\% \;\approx\; \mathbf{103.45\%}.
\end{align*}
The forward clean reproduction of $103.36\%$ is within $\approx 9$~bp
of the paper's quoted value --- attributable to (i) day-count
rounding in $\Adel$ and (ii) the precise repo-rate quote used (paper
gives mid of bid/ask). The forward $\IRR$ is then $y_f$ such that
$\Bond{t_e}{y_f} = \Bf$, which solves to $y_f \approx 2.53\%$
(consistent with the paper).

\subsection{Validation block}

\begin{validation}
This is a \emph{specification} of mathematical objects the combined
study defines and the numerical values they should take on the
canonical case. It says nothing about how those objects are
implemented or invoked. The implementation sequence below is
ordered: each step depends only on the previous ones.

\textbf{Quantities to verify.}
\begin{enumerate}[noitemsep]
\item $\Bond{t}{y}$ in two conventions: \emph{compounded}
      \eqref{eq:cmt-ytm-compounded} and \emph{continuously
      compounded} \eqref{eq:bond-cont}. These are distinct functions;
      only the compounded form defines $\IRR$.
\item $\IRR_t$, the unique zero of $\Bond{t}{y} - P^{mkt}_t$ in the
      compounded convention.
\item $\dD{k}[\Bond{t}{\cdot}]$ for $k=1,2,3$ in closed form
      from \eqref{eq:bond-derivs}.
\item $\PVone(y_f)$ from \eqref{eq:pv01-fwd}.
\item $\Bf$ via the Part-I cash-and-carry formula
      \eqref{eq:bond-forward} (simple compounding) and via the Part-II
      repo-replication formula \eqref{eq:fwd-repo} (continuous
      compounding). They must agree on the canonical case.
\item $y_f$ via the implicit definition \eqref{eq:forward-yield}.
\item The three NPV formulations \eqref{eq:npv-price},
      \eqref{eq:npv-yield}, and the dirty-price derived version.
\item The Taylor remainder $R_1(y)$, satisfying the bounds
      \eqref{eq:rem-bound}.
\item The closed-form greeks $\Delta, \Gamma$ from
      \eqref{eq:delta-P}, \eqref{eq:gamma-full}, and the at-strike
      gamma \eqref{eq:gamma-local}.
\item The sign-threshold formulas \eqref{eq:delta-cond} and
      \eqref{eq:gamma-bound}.
\item The haircut-blended forward via \eqref{eq:fwd-haircut}.
\end{enumerate}

\textbf{Expected numerical outputs on the canonical case.}

\smallskip
\begin{tabular}{@{}p{7cm}rl@{}}
\toprule
Quantity                                        & Value          & Tolerance \\
\midrule
$\Bf$ (forward clean, Part~I)                   & $103.36\%$     & $\pm 0.10\%$ (day-count, repo quote) \\
Forward dirty                                   & $104.74\%$     & $\pm 0.01\%$ \\
$y_f$ (forward IRR)                             & $2.53\%$       & $\pm 1$~bp \\
$\Drepo(t, t_e)$, simple                        & $0.993888$     & $\pm 10^{-5}$ \\
$\Drepo(t, t_e)$, continuous                    & $0.993869$     & $\pm 10^{-5}$ \\
$|R_1|$ at 10Y, $|y-L| \approx 0.3 L$           & $\sim 0.005\%$ & order-of-mag. \\
$M = \max_I |\dD{2}[\Bond{}{\cdot}]|$ (10Y)     & $\sim 150$     & order-of-mag. \\
$\IRR$ gap, compounded vs continuous (10Y at $y\approx 3\%$) & $\sim 4$~bp & order-of-mag. \\
Delta sign threshold $-\dD{1}/\dD{2}$           & positive       & $> 0$ for any coupon schedule \\
Gamma sign threshold $\dD{1}\dD{2}/((\dD{2})^2 - \dD{1}\dD{3})$ & positive & $> 0$ for any coupon schedule \\
\bottomrule
\end{tabular}
\smallskip

\textbf{Implementation sequence (ordered).}
The natural order to build and test, each step depending only on
earlier ones:
\begin{enumerate}[noitemsep]
\item Implement \texttt{dirty\_price\_compounded(y, schedule, c)} per
      \eqref{eq:cmt-ytm-compounded}. \emph{Check:} at
      $y = 2.717\%$ on the canonical schedule, returns $104.1055\%$.
\item Implement \texttt{dirty\_price\_continuous(y, schedule, c)} per
      \eqref{eq:bond-cont}. \emph{Check:} differs from step~1 by a
      few~bp at typical yields; \emph{do not} use step~2 to solve
      for $\IRR$.
\item Root-find for $\IRR$ in step~1. \emph{Check:} given
      $P^{mkt} = 104.1055\%$, returns $\IRR \approx 2.717\%$;
      substituting back in step~1 returns $104.1055\%$.
\item Closed-form yield derivatives \eqref{eq:bond-derivs} from
      step~2. \emph{Check:} numerical bump of step~2 with
      $\Delta y = 1$~bp recovers $\dD{1}$ from \eqref{eq:bond-derivs}
      to float precision; positivity test
      $\dD{2}[\Bond{}{\cdot}] > 0$ on a grid of yields in $[0, 10\%]$.
\item Forward bond price: implement \eqref{eq:bond-forward} (simple,
      Part~I) and \eqref{eq:fwd-repo} (continuous, Part~II) in
      parallel. \emph{Check:} both return $\Bf \approx 103.45\%$ on
      the canonical case (within $\pm 10$~bp of the paper's
      $103.36\%$). The two formulations must agree with each other
      to $\pm 0.01\%$.
\item $y_f$: invert step~5 via root-find on step~2 at expiry.
      \emph{Check:} returns $y_f \approx 2.53\%$.
\item Three NPV formulations. \emph{Check:} all three agree (after
      conventional translation between locked clean / dirty / yield)
      at the canonical case.
\item Overhedge inequality (Proposition~\ref{prop:overhedge}): for
      any $L$, any $y \in [0\%, 10\%]$, and any positive-coupon
      schedule,
      $\Bond{}{L} + \dD{1}[\Bond{}{\cdot}](y - L) \;\le\; \Bond{}{y}$.
\item Local gamma sign (Proposition~\ref{prop:gamma-sign}): at
      $a = +1$ and $y \le L$, $\Gamma < 0$.
\item Sign-threshold positivity: verify numerically that the
      delta threshold $-\dD{1}/\dD{2} > 0$ and the gamma threshold
      $\dD{1}\dD{2}/((\dD{2})^2 - \dD{1}\dD{3}) > 0$ on the canonical
      case, both at $y = L = 2.717\%$.
\item Repo monotonicity: $\FwdP_t(t_e)$ in \eqref{eq:fwd-repo}
      strictly increasing in $\rrepo$ (when spot price exceeds the
      PV of intermediate coupons --- true on canonical case).
\item Haircut limit: \eqref{eq:fwd-haircut} reduces to
      \eqref{eq:fwd-repo} at $h=0$, to unsecured-funded at $h=1$.
\end{enumerate}

\textbf{Two-curve discounting (Part I).} $\Bf$ uses $\Drepo$; the
final cash payoff uses $\Dtl$ (bond zero-yield curve). These are
different curves. In Pucci (Part~II), the formal arbitrage-free
formula uses a single risk-free $D_{tt_e}$, justified by the short
tenor + US-government issuer. The practitioner two-curve approach
is the more general (and more careful) representation.
\end{validation}

% =====================================================================
\begin{todobox}
\begin{itemize}[noitemsep]
\item Reproduce Pucci's Fig.~1 (T-Lock payoff vs forward overhedge) on
      the 24-Jan-2019 case. The remainder $R_1(y)$ is convex with
      minimum (zero) at $y=L$ and small positive values on either
      side.
\item Reproduce Pucci's Figs.~2 \& 3 (Delta, Gamma plots) to visualise
      the sign-flip at extreme yields predicted by
      Propositions~\ref{prop:delta-sign}--\ref{prop:gamma-sign}.
\item Specialness scenario: re-price with $\rrepo$ shifted by
      $-300$~bp (Pucci, p.~12).
\item Repo-bucket sensitivities: bump $\rrepo$ at each coupon node
      independently, get a repo-curve risk decomposition.
\item Reconcile Part~I's two-curve discounting with Part~II's
      single-curve treatment; identify the parameter regime where
      they coincide.
\item Roll-P\&L estimator: reproduce Pucci's Tables 7--10 on
      Treasury auction history.
\item Independent confirmation of the $9$-bp gap between my
      arithmetic ($\Bf = 103.45\%$) and the paper's quoted
      $\Bf = 103.36\%$. Day-count source most likely. Worth verifying
      against an actual bond pricer.
\end{itemize}
\end{todobox}

\end{document}
